% !PRINTED 79
\weberpage{79}

The quantities so defined
\[
        A_{\alpha_1,\alpha_2,\ldots,\alpha_\nu}^{\beta_1,\beta_2,\ldots,\beta_\nu}
\]
\gesperrt{are called the \(\nu\)-th subdeterminants, or subdeterminants of the \(\nu\)-th order.} They are represented by \((n-\nu)\)-rowed determinants. With respect to these subdeterminants, III implies the theorem:

\gesperrt{IV. The subdeterminant \(A_{\alpha_1,\alpha_2,\ldots,\alpha_\nu}^{\beta_1,\beta_2,\ldots,\beta_\nu}\) changes only its sign if two of its lower or two of its upper indices are interchanged; more generally, it remains unaltered in absolute value if the order of the indices \(\alpha_1,\alpha_2,\ldots,\alpha_\nu\) is replaced by any other order, and changes sign or not according as this permutation is of the second or of the first kind.}

If \(\alpha_1',\alpha_2',\ldots,\alpha_\nu'\) denotes any ordering of \(\alpha_1,\alpha_2,\ldots,\alpha_\nu\), then the determinant \(A\) also contains the complex of terms
\[
        \pm a_{\alpha_1'}^{(\beta_1)}a_{\alpha_2'}^{(\beta_2)}\cdots
        a_{\alpha_\nu'}^{(\beta_\nu)}
        A_{\alpha_1,\alpha_2,\ldots,\alpha_\nu}^{\beta_1,\beta_2,\ldots,\beta_\nu}.
\]
Collecting all these terms gives the complex
\begin{equation}
        A_{\alpha_1,\alpha_2,\ldots,\alpha_\nu}^{\beta_1,\beta_2,\ldots,\beta_\nu}
        \Sigma \pm a_{\alpha_1}^{(\beta_1)}a_{\alpha_2}^{(\beta_2)}\cdots
        a_{\alpha_\nu}^{(\beta_\nu)}.
\tag{6}
\end{equation}
The \(\nu\)-rowed determinant which occurs in (6),
\[
        \Sigma\pm a_{\alpha_1}^{(\beta_1)}a_{\alpha_2}^{(\beta_2)}\cdots a_{\alpha_\nu}^{(\beta_\nu)}=
        \begin{vmatrix}
        a_{\alpha_1}^{(\beta_1)}&a_{\alpha_2}^{(\beta_1)}&\cdots&a_{\alpha_\nu}^{(\beta_1)}\\
        a_{\alpha_1}^{(\beta_2)}&a_{\alpha_2}^{(\beta_2)}&\cdots&a_{\alpha_\nu}^{(\beta_2)}\\
        \vdots&\vdots&&\vdots\\
        a_{\alpha_1}^{(\beta_\nu)}&a_{\alpha_2}^{(\beta_\nu)}&\cdots&a_{\alpha_\nu}^{(\beta_\nu)}
        \end{vmatrix},
\]
is called the subdeterminant complementary to
\(A_{\alpha_1,\alpha_2,\ldots,\alpha_\nu}^{\beta_1,\beta_2,\ldots,\beta_\nu}\), and is denoted by
\(B_{\alpha_1,\alpha_2,\ldots,\alpha_\nu}^{\beta_1,\beta_2,\ldots,\beta_\nu}\).
It contains exactly the rows and columns missing in \(A_{\alpha_1,\alpha_2,\ldots,\alpha_\nu}^{\beta_1,\beta_2,\ldots,\beta_\nu}\), and,
%CONT>>
